%%%%%%%%%%%%%%%%%%%%%%%%%%%%%%%%%%%%%%%%%
% Professional Newsletter Template
% LaTeX Template
% Version 1.0 (09/03/14)
%
% Created by:
% Bob Kerstetter (https://www.tug.org/texshowcase/) and extensively modified by:
% Vel (vel@latextemplates.com)
% 
% This template has been downloaded from:
% http://www.LaTeXTemplates.com
%
% License:
% CC BY-NC-SA 3.0 (http://creativecommons.org/licenses/by-nc-sa/3.0/)
%
%%%%%%%%%%%%%%%%%%%%%%%%%%%%%%%%%%%%%%%%%

\documentclass[9pt]{extarticle} % The default font size is 10pt; 11pt and 12pt are alternatives
\usepackage[procnames]{listings}
\usepackage{color}
\usepackage{makecell}

\definecolor{keywords}{RGB}{255,0,90}
\definecolor{comments}{RGB}{0,0,113}
\definecolor{red}{RGB}{160,0,0}
\definecolor{green}{RGB}{0,150,0}
 
\lstset{language=Python, 
        basicstyle=\ttfamily\small, 
        keywordstyle=\color{keywords},
        commentstyle=\color{comments},
        showstringspaces=false,
        procnamekeys={def,class}}

%%%%%%%%%%%%%%%%%%%%%%%%%%%%%%%%%%%%%%%%%
% Professional Newsletter Template
% Structural Definitions File
% Version 1.0 (09/03/14)
%
% Created by:
% Vel (vel@latextemplates.com)
% 
% This file has been downloaded from:
% http://www.LaTeXTemplates.com
%
% License:
% CC BY-NC-SA 3.0 (http://creativecommons.org/licenses/by-nc-sa/3.0/)
%
%%%%%%%%%%%%%%%%%%%%%%%%%%%%%%%%%%%%%%%%%

%----------------------------------------------------------------------------------------
%	REQUIRED PACKAGES
%----------------------------------------------------------------------------------------

\usepackage{listings}
\usepackage{graphicx} % Required for including images
\usepackage{microtype} % Improved typography
\usepackage{multicol} % Used for the two-column layout of the document
\usepackage{booktabs} % Required for nice horizontal rules in tables
\usepackage{wrapfig} % Required for in-line images
\usepackage{float} % Required for forcing figures not to float with the [H] parameter
\usepackage[utf8]{inputenc}
\usepackage{fancyhdr}

%------------------------------------------------
% Fonts

\usepackage{charter} % Use the Charter font as the main document font
\usepackage{courier} % Use the Courier font for \texttt (monospaced) only
\usepackage[T1]{fontenc} % Use T1 font encoding

%------------------------------------------------
% List Separation

\usepackage{enumitem} % Required to customize the list environments
\setlist{noitemsep,nolistsep} % Remove spacing before, after and within lists for a compact look

%------------------------------------------------
% Figure and Table Caption Styles

\usepackage{caption} % Required for changing caption styles
\captionsetup[table]{labelfont={bf,sf},labelsep=period,justification=justified} % Specify the table caption style
\captionsetup[figure]{labelfont={sf,bf},labelsep=period,justification=justified, font=small} % Specify the figure caption style
\setlength{\abovecaptionskip}{10pt} % Whitespace above captions

%------------------------------------------------
% Spacing Between Paragraphs

\makeatletter
\usepackage{parskip}
\setlength{\parskip}{6pt}
\newcommand{\@minipagerestore}{\setlength{\parskip}{6pt}}
\makeatother

%----------------------------------------------------------------------------------------
%	PAGE MARGINS AND SPACINGS
%----------------------------------------------------------------------------------------

\textwidth = 7 in % Text width
\textheight = 10 in % Text height
\oddsidemargin = -18pt % Left side margin on odd pages
\evensidemargin = -18pt % Left side margin on even pages
\topmargin = -36pt % Top margin
\headheight = 0pt % Remove the header by setting its space to 0
\headsep = 0pt % Remove the space between the header and top of the page
\parskip = 2pt % Space between paragraph
\parindent = 0.0in % Paragraph indentation
\pagestyle{empty} % Disable page numbering

%----------------------------------------------------------------------------------------
%	COLORS
%----------------------------------------------------------------------------------------

\usepackage[dvipsnames,svgnames]{xcolor} % Required to specify custom colors

\definecolor{altncolor}{rgb}{.8,0,0} % Dark red
%\definecolor{altncolor}{rgb}{.2,.4,.8} % Dark blue
%\definecolor{altncolor}{rgb}{.84,.16,.16} % Red

\usepackage[colorlinks=true, linkcolor=altncolor, anchorcolor=altncolor, citecolor=altncolor, filecolor=altncolor, menucolor=altncolor, urlcolor=altncolor]{hyperref} % Use the color defined above for all links

%----------------------------------------------------------------------------------------
%	BOX STYLES
%----------------------------------------------------------------------------------------

\usepackage[framemethod=TikZ]{mdframed}% Required for creating boxes
\mdfdefinestyle{sidebar}{
    linecolor=black, % Outer line color
    outerlinewidth=0.5pt, % Outer line width
    roundcorner=0pt, % Amount of corner rounding
    innertopmargin=10pt, % Top margin
    innerbottommargin=10pt, % Bottom margin
    innerrightmargin=10pt, % Right margin
    innerleftmargin=10pt, % Left margin
    backgroundcolor=white, % Box background color
    frametitlebackgroundcolor=white, % Title background color
    frametitlerule=false, % Title rule - true or false
    frametitlerulecolor=white, % Title rule color
    frametitlerulewidth=0.5pt, % Title rule width
    frametitlefont=\Large, % Title heading font specification
    font=\small
}

\mdfdefinestyle{intextbox}{
    linecolor=black, % Outer line color
    outerlinewidth=0.5pt, % Outer line width
    roundcorner=10pt, % Amount of corner rounding
    innertopmargin=7pt, % Top margin
    innerbottommargin=7pt, % Bottom margin
    innerrightmargin=7pt, % Right margin
    innerleftmargin=7pt, % Left margin
    backgroundcolor=white, % Box background color
    frametitlebackgroundcolor=white, % Title background color
    frametitlerule=false, % Title rule - true or false
    frametitlerulecolor=white, % Title rule color
    frametitlerulewidth=0.5pt, % Title rule width
    frametitlefont=\Large % Title heading font specification
}

%----------------------------------------------------------------------------------------
%	HEADING STYLE
%----------------------------------------------------------------------------------------

\newcommand{\heading}[2]{ % Define the \heading command
\vspace{#2} % White space above the heading
{\begin{center}\Large\textbf{#1}\end{center}} % The heading style
\vspace{#2} % White space below the heading
}

\newcommand{\BackToContents}{\hyperlink{contents}{{\small Back to Contents}}} % Define a command for linking back to the contents of the newsletter % Include the document which specifies all packages and structural customizations for this template

\begin{document}

%--------------------------------------------------------------------------------
% HEADER DETAILS
%--------------------------------------------------------------------------------

\pagestyle{fancy}
\fancyhf{}
\chead{segfault@csh.rit.edu}
\rhead{\today}
\lhead{Volume XLVII Issue \#{9}}
\addtolength\footskip{-15px}
\cfoot{"Henry your door is always open, just like your asshole" - 
Spencer (spencerk)}

%----------------------------------------------------------------------------------------
%	HEADER IMAGE
%----------------------------------------------------------------------------------------

\begin{figure}[H]
\centering\vspace{0.5cm}\includegraphics[width=0.8\linewidth]{segfault.png}
\end{figure}

%--------------------------------------------------------------------------------
% HEADER QUOTE
%--------------------------------------------------------------------------------

\vspace{-15px}
\begin{quote}
\centering
\textbf{\textit{I mean what is really the worst that can happen?}}
\end{quote}
\vspace{10px}

%----------------------------------------------------------------------------------------
%	SIDEBAR - FIRST PAGE
%----------------------------------------------------------------------------------------

\vspace{-0.5cm}\begin{minipage}[t]{.30\linewidth} % Mini page taking up 30% of the actual page
\begin{mdframed}[style=sidebar,frametitle={}] % Sidebar box

%-----------------------------------------------------------

\hypertarget{contents}{\textbf{{\large This week on floor\ldots}}} % \hypertarget provides a label to reference using \hyperlink{label}{link text}
\begin{itemize}
\item \hyperlink{firstnews}{Containers with Dgonyeo}
\end{itemize}

\centerline {\rule{.75\linewidth}{.25pt}} % Horizontal line

%-----------------------------------------------------------

\textbf{Notable Upcoming Events:}
Nothing, go home and spend time with your family, or try to avoid them.
I don't know your family situation.

\centerline {\rule{.75\linewidth}{.25pt}} % Horizontal line
\textbf{Scala} \\
One of the best languages out there right now is called Scala. It is able to
effectively combine object-oriented design with functional programming. Scala
runs on top of the JVM which gives you 100\% binary compatibility with a wide
array of library support and resources.  Look forward to the next major release
2.12!\\

\textbf{Python 3.5} \\
Python 3.5, which was release in the middle of September, has several cool new
features. 

\begin{enumerate}[leftmargin=0cm,itemindent=.5cm]
\item \textbf{Coroutines with async and await syntax (PEP 492)} \\
This greatly improves support for asynchronous programming by adding 
coroutines into the language 
\\
\item \textbf{type hints (PEP 484)} \\
While type hints have been in the language since 3.0, however the semantics of
annotations has been left undefined. This introduces a provisional module to 
provide these standard definitions and tools, along with some conventions for
situations where annotations are not available.

\begin{lstlisting}
def msg(name: str) -> str:
    return 'hi ' + name
\end{lstlisting}

\item \textbf{matrix multiplication (PEP 465)} \\
While this might not be a major feature, it still seems like an interesting
one. PEP 465 adds the infix operator `@` for matrix multiplication. 
This was added to help improve Python for data scientists. 
\end{enumerate}

%-----------------------------------------------------------


%-----------------------------------------------------------

%-----------------------------------------------------------

\end{mdframed}
\end{minipage}\hfill % End the sidebar mini page 
%
%----------------------------------------------------------------------------------------
%	MAIN BODY - FIRST PAGE
%----------------------------------------------------------------------------------------
%
\begin{minipage}[t]{.66\linewidth} % Mini page taking up 66% of the actual page
\vspace{-0.4cm}
\hypertarget{firstnews}{\heading{Containers with Dgonyeo}{6pt}}
	
Containers are a really cool, useful technology that are good for solving a lot
of common problems. One problem though is that they have an absurd amount of 
hype behind them, and a lot of people don't fully understand what they are. 

In order for me to describe what containers are, let's first establish what
virtual machines are, and why they exist. Before any of this container or VM
stuff came about, you would just run applications straight on your server. 
If Sally wanted PHP 5.6 on the server to run her stuff, and Bob wanted PHP 5.1
to run something else, they'd either need to complicate their server by having multiple versions of PHP available, or Bob would need to run his stuff on a different machine.

This is annoying, and it's a problem that VMs can solve. Instead of getting a
whole new computer to run his shitty little website, we have software built to
help you put a secondary OS on the same machine through lies. The new OS
\textit{thinks} it's on its own hardware, but all of the hardware bits that
it is twiddling are secretly just emulated by the original OS. The virtual
machine's kernel goes to send a packet to the network interface? Nah, it's 
just a thing in software that goes to the host OS. The kernel tries to write a
file? More lies. That's not a real hard drive, but just a file on the host OS. 
Now Sally can run her new and shiny PHP app, and Bob can run his old and shitty
PHP app, and each of them get their own web server with only the version of PHP
that they want.

Many companies have built fortunes using lies - err, VMs - to be able to give
each little service / application its own machine. One annoying thing about
this strategy though is it can take a significant amount of work to set up a
VM. If you've ever installed Linux, imagine going through that process manually
for each set of things you want to run. There is tooling around doing this
automagically for you, but if you're not paying Amazon you'll be setting that
up yourself.

In our hypothetical scenario where Bob insists on using a version of PHP from
2005 for some god awful reason, what if we could streamline the process by
using significantly less lies? Ultimately the end goal is just having multiple
web servers with different versions of PHP running on the same machine. Do we
really need to emulate network cards, CPUs, hard drives, and more for each web
server? What if we could only lie about the versions of PHP available on the
system, and let the different web servers use the real block devices and memory
space on the host?

Containers are, in a sense, exactly that. Utilizing some features in the Linux
kernel, we can just lie to a process about the following things:
\begin{itemize}
\item Where `/` is in the filesystem. This means that when a process tries to load
  in `/usr/lib/php5`, what it really gets is
  `/var/lib/rkt/stuff/rootfs/usr/lib/php5`.
\item What other processes are running on the system. This means that the two web
  servers won't be able to see - and therefore interfere - with each other (or
  much else for that matter).
\item The network stack. It can be useful to (optionally) hand each web server a
  separate network stack, that is separate from the host's. This means that
  both web servers can bind to port 80, and when the host gets a connection to
  port 80 it'll just forward the connection to one of the two web servers,
  depending on what application the user is trying to access.
\end{itemize}

So there you have it. Containers are kinda like VMs in that they're also built
with the goal of lying to parts of your computer, but they aim to be easier to
use/setup by lying about significantly less.

%-----------------------------------------------------------

\end{minipage} % End the main body - first page mini page

\end{document} 